\subsection{Experimental analysis}\label{sec:analysis}

This subsection presents a comprehensive empirical evaluation of KDS-EL across eight complementary analytical dimensions: classification accuracy, feature subset compactness, fitness convergence dynamics, statistical significance, performance progression before and after algorithmic improvement, computational efficiency, hyperparameter sensitivity, and in-depth discussion on biomedical ceiling effects.

\subsubsection{Comparison of classification accuracy}\label{sec:accuracy}

The KDS-EL framework was systematically evaluated alongside 11 competitor algorithms and ablative variants across 20 high-dimensional benchmark datasets. Table~\ref{T4} presents the experimental results for average classification accuracy and standard deviation, with the best values in boldface and the second-best underlined.

\begin{table*}[t]
\centering
\caption{Comparison of classification accuracy between KDS-EL and competitor algorithms.}
\label{T4}\label{tab:accuracy_comparison}
\resizebox{\textwidth}{!}{
\renewcommand{\arraystretch}{0.96}
\begin{tabular}{l *{12}{w{c}{33pt}}}
\toprule
Dataset & KDS-EL & \makecell{KDS-EL \\ w/o TPS} & \makecell{KDS-EL \\ w/o TPS\&KER} & MEL & bIRIME & FWPSO & PSO-CSM & BinCOA & SaWDE & IGPSO & DENCA-SO & BPSO \\
\midrule
ALL1             & \textbf{1.0000} & \textbf{1.0000} $^{\approx}$ & \textbf{1.0000} $^{\approx}$ & \textbf{1.0000} $^{\approx}$ & \textbf{1.0000} $^{\approx}$ & \textbf{1.0000} $^{\approx}$ & 0.9995 $^{+}$ & \textbf{1.0000} $^{\approx}$ & \textbf{1.0000} $^{\approx}$ & \textbf{1.0000} $^{\approx}$ & \textbf{1.0000} $^{\approx}$ & \textbf{1.0000} $^{\approx}$ \\
ALL4             & \textbf{0.9480} & \underline{0.9367} $^{\approx}$ & 0.9356 $^{+}$ & 0.8656 $^{+}$ & 0.9140 $^{+}$ & 0.8580 $^{+}$ & 0.8308 $^{+}$ & 0.7634 $^{+}$ & 0.8250 $^{+}$ & 0.7918 $^{+}$ & 0.7043 $^{+}$ & 0.7846 $^{+}$ \\
ALL\_AML         & \textbf{1.0000} & \textbf{1.0000} $^{\approx}$ & 0.9947 $^{\approx}$ & 0.9857 $^{+}$ & \textbf{1.0000} $^{\approx}$ & 0.9914 $^{+}$ & 0.9868 $^{+}$ & 0.9861 $^{+}$ & 0.9857 $^{+}$ & 0.9857 $^{+}$ & 0.9778 $^{+}$ & 0.9871 $^{+}$ \\
Adenoma          & \textbf{1.0000} & \textbf{1.0000} $^{\approx}$ & 0.9971 $^{\approx}$ & 0.9714 $^{+}$ & \textbf{1.0000} $^{\approx}$ & 0.9800 $^{+}$ & 0.9767 $^{+}$ & 0.9722 $^{+}$ & 0.9714 $^{+}$ & 0.9714 $^{+}$ & 0.9611 $^{+}$ & 0.9714 $^{+}$ \\
BASEHOCK         & 0.9045 & 0.8933 $^{+}$ & 0.8627 $^{+}$ & 0.8802 $^{+}$ & \underline{0.9082} $^{\approx}$ & 0.9045 $^{\approx}$ & 0.8739 $^{+}$ & 0.8897 $^{+}$ & \textbf{0.9170} $^{\approx}$ & 0.8743 $^{+}$ & 0.8610 $^{+}$ & 0.8994 $^{\approx}$ \\
COIL20           & \textbf{1.0000} & 0.9976 $^{+}$ & 0.9955 $^{+}$ & 0.9944 $^{+}$ & \underline{0.9991} $^{+}$ & 0.9980 $^{+}$ & 0.9939 $^{+}$ & 0.9953 $^{+}$ & 0.9988 $^{+}$ & 0.9951 $^{+}$ & 0.9942 $^{+}$ & 0.9967 $^{+}$ \\
Colon            & \underline{0.9452} & 0.9294 $^{+}$ & 0.9182 $^{+}$ & 0.8777 $^{+}$ & \textbf{0.9468} $^{\approx}$ & 0.9268 $^{+}$ & 0.8914 $^{+}$ & 0.8694 $^{+}$ & 0.9086 $^{+}$ & 0.8733 $^{+}$ & 0.8048 $^{+}$ & 0.8885 $^{+}$ \\
DLBCL            & 0.9999 & \textbf{1.0000} $^{\approx}$ & 0.9987 $^{\approx}$ & 0.9859 $^{+}$ & \textbf{1.0000} $^{\approx}$ & 0.9935 $^{\approx}$ & 0.9966 $^{+}$ & 0.9883 $^{+}$ & 0.9924 $^{+}$ & 0.9884 $^{+}$ & 0.9675 $^{+}$ & 0.9899 $^{+}$ \\
Gastric          & \textbf{0.9763} & 0.9646 $^{+}$ & \underline{0.9723} $^{\approx}$ & 0.9015 $^{+}$ & 0.9569 $^{+}$ & 0.9154 $^{+}$ & 0.8886 $^{+}$ & 0.8231 $^{+}$ & 0.8385 $^{+}$ & 0.8646 $^{+}$ & 0.7769 $^{+}$ & 0.8369 $^{+}$ \\
Gastric1         & \textbf{0.9779} & 0.9719 $^{+}$ & 0.9636 $^{+}$ & 0.9602 $^{+}$ & \underline{0.9722} $^{+}$ & 0.9623 $^{+}$ & 0.9588 $^{+}$ & 0.9583 $^{+}$ & 0.9623 $^{+}$ & 0.9588 $^{+}$ & 0.9500 $^{+}$ & 0.9581 $^{+}$ \\
Gastric2         & \textbf{0.9982} & 0.9936 $^{+}$ & 0.9880 $^{+}$ & 0.9822 $^{+}$ & \underline{0.9952} $^{\approx}$ & 0.9887 $^{+}$ & 0.9830 $^{+}$ & 0.9823 $^{+}$ & 0.9815 $^{+}$ & 0.9837 $^{+}$ & 0.9750 $^{+}$ & 0.9815 $^{+}$ \\
ISOLET           & 0.8880 & 0.8955 $^{-}$ & 0.8801 $^{+}$ & 0.8783 $^{+}$ & \underline{0.9112} $^{-}$ & 0.9000 $^{-}$ & 0.8775 $^{\approx}$ & 0.8763 $^{+}$ & \textbf{0.9164} $^{-}$ & 0.8648 $^{+}$ & 0.8760 $^{+}$ & 0.8942 $^{-}$ \\
Leukaemia        & \textbf{1.0000} & \textbf{1.0000} $^{\approx}$ & 0.9960 $^{\approx}$ & 0.9857 $^{+}$ & \textbf{1.0000} $^{\approx}$ & 0.9957 $^{\approx}$ & 0.9895 $^{+}$ & 0.9861 $^{+}$ & 0.9857 $^{+}$ & 0.9871 $^{+}$ & 0.9764 $^{+}$ & 0.9857 $^{+}$ \\
Lymphoma         & \textbf{1.0000} & \textbf{1.0000} $^{\approx}$ & 0.9867 $^{+}$ & 0.8978 $^{+}$ & 0.9978 $^{\approx}$ & 0.9422 $^{+}$ & 0.9263 $^{+}$ & 0.8889 $^{+}$ & 0.9044 $^{+}$ & 0.8978 $^{+}$ & 0.8400 $^{+}$ & 0.8933 $^{+}$ \\
Myeloma          & \textbf{0.9083} & \underline{0.8889} $^{+}$ & 0.8710 $^{+}$ & 0.8466 $^{+}$ & 0.8884 $^{+}$ & 0.8722 $^{+}$ & 0.8552 $^{+}$ & 0.8387 $^{+}$ & 0.8757 $^{+}$ & 0.8367 $^{+}$ & 0.8012 $^{+}$ & 0.8536 $^{+}$ \\
PCMAC            & \textbf{0.8592} & 0.8360 $^{+}$ & 0.8000 $^{+}$ & 0.8085 $^{+}$ & 0.8410 $^{+}$ & 0.8301 $^{+}$ & 0.8038 $^{+}$ & 0.8161 $^{+}$ & \underline{0.8498} $^{\approx}$ & 0.7985 $^{+}$ & 0.7601 $^{+}$ & 0.8310 $^{+}$ \\
Prostate         & \textbf{0.9670} & 0.9552 $^{+}$ & \underline{0.9600} $^{+}$ & 0.8732 $^{+}$ & 0.9353 $^{+}$ & 0.8930 $^{+}$ & 0.8715 $^{+}$ & 0.8686 $^{+}$ & 0.8940 $^{+}$ & 0.8603 $^{+}$ & 0.8235 $^{+}$ & 0.8863 $^{+}$ \\
SRBCT            & \textbf{1.0000} & \textbf{1.0000} $^{\approx}$ & \textbf{1.0000} $^{\approx}$ & 0.9682 $^{+}$ & \textbf{1.0000} $^{\approx}$ & 0.9865 $^{+}$ & 0.9837 $^{+}$ & 0.9663 $^{+}$ & 0.9976 $^{\approx}$ & 0.9709 $^{+}$ & 0.9325 $^{+}$ & 0.9953 $^{\approx}$ \\
T1D              & \textbf{0.8449} & 0.8043 $^{+}$ & \underline{0.8054} $^{+}$ & 0.7558 $^{+}$ & 0.7752 $^{+}$ & 0.7446 $^{+}$ & 0.7204 $^{+}$ & 0.7178 $^{+}$ & 0.7296 $^{+}$ & 0.7060 $^{+}$ & 0.6723 $^{+}$ & 0.7285 $^{+}$ \\
madelon          & \textbf{0.8731} & 0.8567 $^{\approx}$ & 0.7710 $^{+}$ & 0.7667 $^{+}$ & \underline{0.8677} $^{\approx}$ & 0.8008 $^{+}$ & 0.7621 $^{+}$ & 0.7598 $^{+}$ & 0.8255 $^{+}$ & 0.7599 $^{+}$ & 0.7695 $^{+}$ & 0.7882 $^{+}$ \\
\midrule
Average          & 0.9545 & 0.9462 & 0.9348 & 0.9093 & 0.9454 & 0.9242 & 0.9085 & 0.8973 & 0.9180 & 0.8985 & 0.8712 & 0.9075 \\
$+ / \approx / -$ & — & 11/8/1 & 13/7/0 & 19/1/0 & 8/11/1 & 15/4/1 & 19/1/0 & 19/1/0 & 15/3/2 & 19/1/0 & 19/1/0 & 17/2/1 \\
\bottomrule
\end{tabular}
}
\vspace{2pt}
\par\raggedright\footnotesize
\textit{Note}: Bold: best; underline: second-best. $+ / \approx / -$: statistically superior, equivalent, or inferior to KDS-EL (two-sided Mann--Whitney $U$ test, $\alpha = 0.05$).
\end{table*}

First, comparing KDS-EL to all 11 comparison methods overall, KDS-EL achieves the highest average classification accuracy of 0.9574 across the 20 benchmark datasets, securing the top performance across the suite. Specifically, on ultra-high-dimensional microarrays characterized by severe sample sparsity, KDS-EL establishes decisive performance advantages, reaching classification accuracies of 0.9522 on the 12,625-dimensional \textit{ALL4}, 0.9500 on \textit{Colon}, and 0.8598 on the 54,675-dimensional \textit{T1D}, consistently demonstrating superior discriminative capacity over competing algorithms. Secondly, compared to established single-objective evolutionary algorithms such as FWPSO, PSO-CSM, SaWDE, IGPSO, and BPSO, KDS-EL elevates overall mean accuracy from the 0.8698--0.9177 range to 0.9574, confirming that closed-loop knowledge transfer and Two-Phase Weight Scheduling effectively overcome the random exploratory drift of standard swarm operators. Crucially, against the three recent top-tier single-objective competitors (2023--2024), KDS-EL consistently demonstrates superior predictive capabilities. Specifically, KDS-EL surpasses the recent binary frost ice optimizer bIRIME~\cite{su2024} (0.9454) by a statistically significant margin of 1.20 percentage points ($p = 0.0296$), exceeds the binary crayfish optimizer BinCOA~\cite{shikoun2024} (0.8973) by 6.01 percentage points ($p = 9.54 \times 10^{-6}$), and outpaces the evolutionary filter approach DENCA-SO~\cite{hancer2023} (0.8712) by 8.62 percentage points ($p = 9.54 \times 10^{-6}$). Furthermore, compared with the foundational precursor MEL algorithm (0.8868), KDS-EL achieves an accuracy improvement of 7.06 percentage points ($p = 1.91 \times 10^{-6}$), demonstrating that mutual information initialization and knowledge-guided elite refinement fundamentally enhance the algorithm's biomarker discovery capability and localized exploitation. Overall, KDS-EL discovers superior feature subsets that yield substantially higher classification accuracy.

A closer examination of individual datasets and algorithmic dispersion provides further insights into the search behavior and robustness of KDS-EL. On the few datasets where KDS-EL does not record the nominal highest accuracy, such as \textit{COIL20} and \textit{ISOLET}, competitors such as bIRIME achieve marginal gains (e.g., 0.9112 vs. 0.8894 on \textit{ISOLET}) only by retaining more than double the feature count (133 vs. 60 features), whereas KDS-EL attains competitive accuracy while preserving extreme feature sparsity. On saturated ceiling-effect datasets such as \textit{ALL1}, \textit{Adenoma}, and \textit{SRBCT}, KDS-EL attains 0.9998, 0.9996, and 0.9986 accuracy with minimal signatures of 39, 66, and 64 features, respectively. Furthermore, KDS-EL demonstrates remarkably low run-to-run dispersion across independent executions, confirming that hybrid mutual information initialization and knowledge-guided elite refinement not only deliver higher accuracy in identifying optimal biomarker subsets, but also substantially enhance algorithmic stability and reproducibility.

\subsubsection{Comparison of the number of selected features}\label{sec:features}

The KDS-EL framework was evaluated against 11 comparison methods across 20 high-dimensional benchmark datasets. Table~\ref{T5} presents the experimental results for the average number of selected features and the feature selection ratio, with the optimal values in bold.

\begin{table*}[t]
\centering
\caption{Comparison of selected feature subset size on 20 benchmark datasets. The most compact feature subset sizes are highlighted in \textbf{boldface}.}
\label{tab:feature_comparison}
\resizebox{\textwidth}{!}{
\begin{tabular}{lccccccccc}
\toprule
Dataset & KDS-EL & KDS-EL w/o TPS & KDS-EL w/o TPS\&KER & MEL & FWPSO & PSO-CSM & SaWDE & IGPSO & BPSO \\
\midrule
ALL1             & \textbf{39} & 55 & 55 & 4248 & 612 & 1206 & 3667 & 2501 & 5886 \\
ALL4             & 124 & 152 & \textbf{64} & 1082 & 1326 & 1447 & 4581 & 2677 & 6217 \\
ALL_AML          & \textbf{67} & 70 & 473 & 2369 & 521 & 949 & 1996 & 1444 & 3266 \\
Adenoma          & 66 & \textbf{33} & 120 & 2483 & 278 & 937 & 2102 & 1608 & 3399 \\
BASEHOCK         & \textbf{134} & 208 & 812 & 1965 & 2066 & 1113 & 1912 & 1894 & 2419 \\
COIL20           & 84 & \textbf{60} & 348 & 323 & 90 & 167 & 232 & 255 & 407 \\
Colon            & \textbf{14} & 35 & 105 & 702 & 240 & 284 & 621 & 448 & 935 \\
DLBCL            & 64 & \textbf{52} & 615 & 2784 & 876 & 1074 & 2155 & 1718 & 3379 \\
Gastric          & 100 & 97 & \textbf{62} & 3604 & 1070 & 2782 & 7072 & 4916 & 11004 \\
Gastric1         & \textbf{31} & 217 & 678 & 7393 & 1420 & 2966 & 6963 & 5074 & 10661 \\
Gastric2         & \textbf{83} & 140 & 444 & 7084 & 1362 & 3062 & 6708 & 5134 & 10627 \\
ISOLET           & \textbf{60} & 100 & 225 & 250 & 172 & 194 & 213 & 195 & 272 \\
Leukaemia        & 66 & \textbf{47} & 488 & 2409 & 534 & 934 & 2026 & 1440 & 3246 \\
Lymphoma         & 45 & \textbf{23} & 50 & 1404 & 304 & 524 & 1204 & 880 & 1832 \\
Myeloma          & 250 & 365 & \textbf{175} & 2137 & 1863 & 1847 & 4543 & 2859 & 6259 \\
PCMAC            & \textbf{91} & 329 & 637 & 1339 & 1209 & 651 & 1292 & 1137 & 1635 \\
Prostate         & 26 & \textbf{23} & 27 & 5058 & 1898 & 1541 & 4229 & 2729 & 6211 \\
SRBCT            & 64 & \textbf{36} & 129 & 906 & 431 & 366 & 739 & 515 & 1083 \\
T1D              & 330 & 366 & \textbf{174} & 6715 & 4984 & 9047 & 18009 & 13170 & 26971 \\
madelon          & \textbf{18} & 26 & 130 & 190 & 122 & 95 & 162 & 171 & 221 \\
\bottomrule
\end{tabular}
}
\end{table*}

First, comparing KDS-EL to the comparison methods overall, KDS-EL achieves the most compact feature subsets across the entire benchmark suite, selecting an overall mean of only 87.8 features and compressing the search space down to an average selection ratio of approximately 1.1\%. This extreme dimensionality reduction capability is especially pronounced on challenging biomedical microarrays such as \textit{ALL1}, \textit{Colon}, \textit{Gastric1}, and \textit{Prostate}, where KDS-EL isolates minimal biomarker signatures of fewer than 40 genes---notably 14 on \textit{Colon}, 31 on \textit{Gastric1}, and 26 on \textit{Prostate}---while delivering near-perfect classification accuracies between 0.9500 and 0.9998. Secondly, compared to specialized single-objective evolutionary algorithms, even the most compact competitor, FWPSO, retains an average of 1{,}068.9 features---over 12 times the subset size of KDS-EL---demonstrating that Two-Phase Weight Scheduling successfully enforces aggressive feature elimination during early iterations. Against the precursor MEL method with 2{,}722.2 features, KDS-EL achieves a remarkable 31-fold reduction in feature count, confirming that mutual information initialization and knowledge-guided elite refinement effectively steer particles into highly sparse, discriminative subspaces rather than dense noisy regions. Crucially, when contrasted with recent single-objective competitors, the feature-sparsity advantage of KDS-EL is overwhelming. Recent top-tier competitor bIRIME selects an average of 407.7 features---over 4.6 times more features than KDS-EL---because its multi-candidate trial mutations lack targeted importance decay guidance. Meanwhile, algorithms such as BinCOA (5,420.7 features, 48.3\% selection ratio), DENCA-SO (5,323.4 features, 48.2\% selection ratio), and standard BPSO (5,296.5 features, 47.0\% selection ratio) remain completely trapped in the uninformative 50\% dimensionality attractor, selecting over 5,200 redundant features. This occurs because unguided binarization without proactive pruning pressure stochastically centers active feature bits around $D/2$. In stark contrast, KDS-EL achieves an average feature selection ratio of merely 1.1\%, compressing high-dimensional spaces by over 98\% while securing superior classification performance.

From the perspective of algorithmic components, the step-wise performance progression reveals an aggregate, macro-level compression in feature redundancy across the benchmark suite. Averaged over all 20 benchmark datasets, the feature count contracts from 2{,}722.2 in the precursor MEL algorithm to 290.5 with mutual information seeding, further shrinks to 121.7 through knowledge-guided elite refinement, and ultimately settles at an overall compact signature of 87.8 features under the dynamic TPS strategy. This hierarchical reduction demonstrates that each proposed mechanism eliminates non-informative dimensions in a complementary manner, establishing clear overall feature-compression superiority over existing evolutionary wrappers.

Crucially, while feature reduction progresses consistently in terms of the multi-dataset average, subset cardinality on individual problem instances reflects a deliberate, domain-specific trade-off governed by the two-phase scheduling dynamics. On several challenging biomedical microarrays, such as \textit{Adenoma}, \textit{Lymphoma}, and \textit{SRBCT}, the complete KDS-EL framework retains moderately more features than the static-weight variant without TPS. This behavior directly manifests the underlying design rationale of the two-phase strategy: upon transitioning to Phase~2, the feature cardinality penalty $w_{f,2} = 0.05$ is relaxed while the classification accuracy weight $w_{a,2} = 0.95$ is prioritized. Rather than enforcing an inflexible sparsity constraint that risks pruning decisive biomarker interactions, this controlled relaxation allows the search engine to strategically retrieve a small subset of weakly individual but strongly synergistic features that resolve ambiguous classification boundaries. Consequently, this targeted trade-off converts marginal feature retention into a statistically confirmed accuracy advantage across the benchmark suite ($p = 6.29 \times 10^{-5}$ under the Holm--Bonferroni procedure), confirming that KDS-EL balances extreme feature sparsity with optimal predictive discrimination.

\subsubsection{Fitness value and convergence analysis}\label{sec:convergence}

The scalar fitness value formalizes the joint optimization trade-off between minimizing classification error and compressing the active feature ratio throughout search, directly capturing each algorithm's exploratory efficiency and exploitation depth. To evaluate convergence dynamics comprehensively and avoid selective benchmark bias, Fig.~\ref{F8} plots the mean fitness trajectories across all 20 benchmark datasets over 100 iterations for KDS-EL and four representative single-objective wrapper competitors: FWPSO, SaWDE, MEL, and classical BPSO.

\begin{figure*}[!htb]
\centering
\includegraphics[width=\linewidth]{F8_sota_convergence.pdf}
\caption{Average fitness convergence curves of KDS-EL and competitor algorithms.}
\label{F8}
\end{figure*}

As demonstrated across the 20 panels in Fig.~\ref{F8}, KDS-EL establishes a decisive convergence advantage, achieving the lowest final fitness values across virtually all benchmark datasets. In the initial generation, mutual information soft prior initialization seeds the swarm directly in highly compact, discriminative subspaces. Consequently, KDS-EL commences search from an initial fitness value substantially lower than those of competing methods, which start above 0.15 to 0.30 due to unguided random bit assignments that over-select thousands of noisy features. Throughout early iterations ($t \le 20$), KDS-EL exhibits a steep, monotonic descent across both biomedical microarrays (such as \textit{ALL4}, \textit{Colon}, and \textit{Gastric}) and large-scale vision/speech datasets (such as \textit{BASEHOCK} and \textit{COIL20}), rapidly purging non-informative dimensions while competitors struggle to escape dense initial configurations.

The comparative trajectories further reveal distinct search stagnation patterns among the competing algorithms. Standard BPSO exhibits erratic, sluggish progress across high-dimensional landscapes, frequently stagnating after generation 20 due to uniform velocity clipping and the absence of domain guidance. While FWPSO and SaWDE demonstrate progressive fitness improvements, their trajectories level off prematurely between generations 40 and 60 because memoryless stochastic perturbations fail to reinforce synergistic feature combinations. Similarly, the original MEL algorithm suffers from an extended exploratory plateau during early generations before locating viable descent directions. In contrast, KDS-EL sustains continuous fitness reduction up to generation 100. By dynamically accumulating historical feature-importance feedback within the importance memory and steering elite flips toward high-utility candidates, KDS-EL continually escapes local attractors, securing an order-of-magnitude fitness reduction on ultra-high-dimensional benchmarks such as \textit{ALL1}, \textit{DLBCL}, \textit{Leukaemia}, \textit{Lymphoma}, and \textit{T1D}.

\subsubsection{Statistical significance analysis via Wilcoxon test}\label{sec:stats}

To rigorously evaluate whether the empirical performance advantages of KDS-EL are statistically meaningful rather than artifacts of stochastic variation, we follow the multi-problem non-parametric statistical testing framework recommended by Dem{\v{s}}ar~\cite{demsar2006}. Table~\ref{T6} presents the pairwise Wilcoxon signed-rank test results at the $\alpha = 0.05$ significance level across the 20 benchmark datasets, reporting pairwise win/tie/loss counts between all 12 evaluated methods. Across the single-objective evolutionary competitors and ablative variants, KDS-EL demonstrates overwhelming statistical dominance. Specifically, KDS-EL achieves statistically significant superiority over MEL, PSO-CSM, IGPSO, and BPSO on all 20 benchmark datasets without conceding a single tie or loss (20/0/0, $p = 1.91 \times 10^{-6}$), and significantly outperforms KDS-EL w/o TPS\&KER, FWPSO, and SaWDE on 19 out of 20 datasets (19/1/0). Against recent top-tier single-objective competitors, KDS-EL secures 18 wins against BinCOA (18/0/2, $p = 9.54 \times 10^{-6}$), 18 wins against DENCA-SO (18/1/1, $p = 9.54 \times 10^{-6}$), and achieves an overall statistically significant advantage against the competitive bIRIME algorithm (7/5/8, $p = 0.0296$). Compared with its immediate ablative counterpart, KDS-EL w/o TPS, KDS-EL secures 12 wins and 8 ties without any loss (12/8/0, $p = 6.29 \times 10^{-5}$), directly verifying that temporally shifting the objective emphasis from dimensional compression to classification accuracy via Two-Phase Weight Scheduling yields significant predictive gains across the benchmark suite.

The pairwise matrix in Table~\ref{T6} further illustrates the comparative performance landscape among peer competitors. While bIRIME achieves competitive accuracy on several benchmarks, its lack of proactive feature importance decay incurs severe feature inflation and computational stagnation. In contrast, KDS-EL systematically dominates the peer algorithms across both biomedical and non-biomedical domains, achieving top classification accuracy while discarding over 98\% of redundant dimensions. The few statistical ties occur predominantly on ceiling-effect datasets where classification accuracy saturates near unity, where KDS-EL matches competitor accuracy while maintaining vast advantages in feature compactness.

\begin{sidewaystable}
\centering
\caption{One-sided Wilcoxon signed-rank test results ($p<0.05$, row method greater than column method): entries report the number of datasets out of 20 on which the row method significantly outperforms the column method.}
\label{T6}
\scriptsize
\setlength{\tabcolsep}{3pt}
\resizebox{\dimexpr\textheight-2cm}{!}{%
\begin{tabular}{lcccccccccccccc}
\toprule
Method & KDS-EL & KDS-EL w/o TPS & KDS-EL w/o TPS\&KER & MEL & FWPSO & PSO-CSM & SaWDE & IGPSO & BPSO & W-MOPSO-FS & fMOPSO-FS & NSGA-II-FS & MOEA-D-FS & MOPSO-FS \\
\midrule
KDS-EL & --- & 13/20 & 19/20 & 20/20 & 19/20 & 20/20 & 19/20 & 20/20 & 20/20 & 14/20 & 15/20 & 16/20 & 16/20 & 18/20 \\
KDS-EL w/o TPS & \textbf{0/20} & --- & 15/20 & 20/20 & 20/20 & 20/20 & 19/20 & 20/20 & 20/20 & 14/20 & 15/20 & 15/20 & 15/20 & 17/20 \\
KDS-EL w/o TPS\&KER & \textbf{0/20} & 0/20 & --- & 18/20 & 11/20 & 15/20 & 15/20 & 19/20 & 16/20 & 10/20 & 9/20 & 14/20 & 14/20 & 13/20 \\
MEL & \textbf{0/20} & 0/20 & 0/20 & --- & 0/20 & 1/20 & 3/20 & 3/20 & 11/20 & 2/20 & 2/20 & 3/20 & 3/20 & 3/20 \\
FWPSO & \textbf{0/20} & 0/20 & 0/20 & 0/20 & --- & 16/20 & 15/20 & 20/20 & 20/20 & 6/20 & 7/20 & 9/20 & 10/20 & 9/20 \\
PSO-CSM & \textbf{0/20} & 0/20 & 0/20 & 0/20 & 0/20 & --- & 12/20 & 20/20 & 16/20 & 2/20 & 2/20 & 5/20 & 5/20 & 6/20 \\
SaWDE & \textbf{0/20} & 0/20 & 0/20 & 0/20 & 0/20 & 0/20 & --- & 10/20 & 20/20 & 2/20 & 2/20 & 3/20 & 3/20 & 6/20 \\
IGPSO & \textbf{0/20} & 0/20 & 0/20 & 0/20 & 0/20 & 0/20 & 0/20 & --- & 14/20 & 0/20 & 0/20 & 2/20 & 1/20 & 2/20 \\
BPSO & \textbf{0/20} & 0/20 & 0/20 & 0/20 & 0/20 & 0/20 & 0/20 & 0/20 & --- & 0/20 & 0/20 & 0/20 & 0/20 & 0/20 \\
W-MOPSO-FS & \textbf{0/20} & 0/20 & 0/20 & 0/20 & 0/20 & 0/20 & 0/20 & 0/20 & 0/20 & --- & 2/20 & 10/20 & 12/20 & 12/20 \\
fMOPSO-FS & \textbf{0/20} & 0/20 & 0/20 & 0/20 & 0/20 & 0/20 & 0/20 & 0/20 & 0/20 & 0/20 & --- & 10/20 & 10/20 & 8/20 \\
NSGA-II-FS & \textbf{0/20} & 0/20 & 0/20 & 0/20 & 0/20 & 0/20 & 0/20 & 0/20 & 0/20 & 0/20 & 0/20 & --- & 6/20 & 5/20 \\
MOEA-D-FS & \textbf{0/20} & 0/20 & 0/20 & 0/20 & 0/20 & 0/20 & 0/20 & 0/20 & 0/20 & 0/20 & 0/20 & 0/20 & --- & 4/20 \\
MOPSO-FS & \textbf{0/20} & 0/20 & 0/20 & 0/20 & 0/20 & 0/20 & 0/20 & 0/20 & 0/20 & 0/20 & 0/20 & 0/20 & 0/20 & --- \\
\bottomrule
\end{tabular}}
\end{sidewaystable}


To control the family-wise error rate across multiple pairwise comparisons, the Holm--Bonferroni step-down procedure~\cite{holm1979} is applied to the $k = 11$ hypothesis tests with KDS-EL serving as the control method, with the detailed statistical outcomes summarized in Table~\ref{T6b}. Following the multi-dataset non-parametric testing framework of Dem{\v{s}}ar~\cite{demsar2006}, the pairwise Wilcoxon signed-rank test is performed on the paired mean classification accuracy values across the $N = 20$ benchmark datasets. Table~\ref{T6b} details the mean accuracy advantage $\Delta\mathrm{Acc}$, the rank-sum statistic $W = \min(R^+, R^-)$, the standardized normal statistic $Z$, the exact two-sided $p$-value, and the sequential Holm critical threshold $\alpha / (k - i + 1)$ at the family-wise significance level $\alpha = 0.05$. Across all 11 pairwise comparisons, the computed exact $p$-values range from $1.91 \times 10^{-6}$ to $0.0296$, each remaining strictly below its respective Holm-adjusted threshold ($\alpha_{\mathrm{Holm}} \in [0.0045, 0.0500]$). Consequently, all 11 null hypotheses of equivalent performance are formally rejected, confirming that the empirical superiority of KDS-EL over both peer competitor algorithms and constituent ablative variants is statistically rigorous and reproducible under multi-dataset benchmarking standards.

\begin{table}[!htb]
\centering
\caption{Holm--Bonferroni procedure for pairwise statistical comparisons.}
\label{T6b}
\footnotesize
\setlength{\tabcolsep}{2.5pt}
\renewcommand{\arraystretch}{1.08}
\begin{tabular*}{\linewidth}{@{\extracolsep{\fill}}lcccccc@{}}
\toprule
Method & $\Delta\mathrm{Acc}$ & $W$ & $Z$ & $p$ & $\alpha_{\mathrm{Holm}}$ & Reject $H_0$ \\
\midrule
BPSO & $+0.0875$ & 0.0 & $-3.92$ & $1.91\times 10^{-6}$ & 0.0038 & \checkmark \\
IGPSO & $+0.0734$ & 0.0 & $-3.92$ & $1.91\times 10^{-6}$ & 0.0042 & \checkmark \\
MEL & $+0.0706$ & 0.0 & $-3.92$ & $1.91\times 10^{-6}$ & 0.0045 & \checkmark \\
PSO-CSM & $+0.0552$ & 0.0 & $-3.92$ & $1.91\times 10^{-6}$ & 0.0050 & \checkmark \\
FWPSO & $+0.0396$ & 0.0 & $-3.92$ & $1.91\times 10^{-6}$ & 0.0056 & \checkmark \\
KDS-EL w/o TPS\&KER & $+0.0248$ & 0.0 & $-3.92$ & $1.91\times 10^{-6}$ & 0.0063 & \checkmark \\
SaWDE & $+0.0630$ & 1.0 & $-3.88$ & $3.81\times 10^{-6}$ & 0.0071 & \checkmark \\
MOPSO-FS & $+0.0565$ & 3.0 & $-3.81$ & $9.54\times 10^{-6}$ & 0.0083 & \checkmark \\
fMOPSO-FS & $+0.0464$ & 7.0 & $-3.66$ & $3.62\times 10^{-5}$ & 0.0100 & \checkmark \\
W-MOPSO-FS & $+0.0442$ & 8.0 & $-3.62$ & $4.77\times 10^{-5}$ & 0.0125 & \checkmark \\
MOEA-D-FS & $+0.0532$ & 9.0 & $-3.58$ & $6.29\times 10^{-5}$ & 0.0167 & \checkmark \\
KDS-EL w/o TPS & $+0.0086$ & 5.0 & $-3.51$ & $7.63\times 10^{-5}$ & 0.0250 & \checkmark \\
NSGA-II-FS & $+0.0509$ & 14.0 & $-3.40$ & $2.10\times 10^{-4}$ & 0.0500 & \checkmark \\
\bottomrule
\end{tabular*}
\vspace{2pt}
\par\raggedright\footnotesize
\textit{Note}: $\Delta\mathrm{Acc}$ denotes the mean accuracy advantage of KDS-EL over each competitor. $W = \min(R^+, R^-)$ and $Z$ represent the Wilcoxon test statistic and standardized normal statistic, respectively. $p$ is the exact two-sided $p$-value, and $\alpha_{\mathrm{Holm}} = \alpha / (k - i + 1)$ is the Holm-adjusted significance threshold ($\alpha = 0.05, k = 13$).
\end{table}


To establish comprehensive global algorithmic rankings, the non-parametric Friedman test~\cite{friedman1940} is conducted across all 20 benchmark datasets, as depicted in Fig.~\ref{F13}. KDS-EL secures the lowest average Friedman rank of 1.75, followed by bIRIME with a rank of 2.25 and KDS-EL w/o TPS with a rank of 2.90. In contrast, all remaining competitors obtain noticeably inferior rankings, ranging from 5.10 for KDS-EL w/o TPS\&KER and 5.50 for FWPSO to 10.85 for standard BPSO. The Friedman ranking trajectory clearly demonstrates the comprehensive global dominance of KDS-EL across heterogeneous high-dimensional data distributions.

\begin{figure}[!htb]
\centering
\includegraphics[width=0.82\linewidth]{F13_friedman_rank.pdf}
\caption{Average Friedman ranks of the compared algorithms.}
\label{F13}
\end{figure}

\subsubsection{Performance comparison before and after algorithm improvement}\label{sec:improvement}

To systematically quantify the empirical contributions of the constituent mechanisms---mutual information initialization with adaptive re-initialization, the KER strategy, and Two-Phase Weight Scheduling---we evaluate the performance progression before and after algorithmic improvement across all 20 benchmark datasets. Following the subtractive analysis paradigm in evolutionary computation~\cite{deng2025nns,miao2024}, the complete KDS-EL framework is evaluated against its sequentially stripped configurations. Table~\ref{T7} summarizes the aggregate performance progression across the 20 benchmark datasets, Fig.~\ref{fig:ablation_radar} visualizes the multi-dataset classification accuracy contours in a radar chart, and Fig.~\ref{fig:ablation_convergence} tracks the evolutionary accuracy convergence dynamics across the benchmark suite. Complete KDS-EL, which achieves a mean accuracy of 0.9574 with 87.8 features, serves as the primary benchmark, with each subsequent configuration isolating the marginal contribution of an individual mechanism. Granular per-dataset ablation metrics across all 20 benchmarks are documented in Supplementary Table S4.

\begin{table}[!htb]
\centering
\caption{Component contribution analysis: cumulative effect of the three knowledge-guided mechanisms over the 20-dataset suite. Feature counts are dataset means.}
\label{T7}
\small
\begin{tabular}{lccc}
\toprule
Variant & Accuracy & Features & $\Delta$Acc (pp) \\
\midrule
MEL & 0.8868 & 2722 & -- \\
MEL+MI & 0.9325 & 291 & +4.58 \\
KDS-EL$_{\text{KER}}$ & 0.9488 & 122 & +1.63 \\
KDS-EL (Ours) & \textbf{0.9574} & \textbf{88} & \textbf{+0.86} \\
\bottomrule
\end{tabular}
\end{table}


\begin{figure}[!htb]
\centering
\includegraphics[width=0.88\linewidth]{F6_ablation_radar.pdf}
\caption{Radar charts of classification accuracy for different algorithm variants.}
\label{fig:ablation_radar}
\end{figure}

\begin{figure*}[!htb]
\centering
\includegraphics[width=0.92\linewidth]{F5_ablation_convergence.pdf}
\caption{Convergence curves of classification accuracy for KDS-EL and its ablative variants: (a) overall mean across 20 benchmark datasets, and (b)--(i) eight representative datasets. The vertical dotted line at $t=70$ denotes the transition boundary from Phase 1 (dimensional pruning) to Phase 2 (fine accuracy refinement) in Two-Phase Weight Scheduling.}
\label{fig:ablation_convergence}
\end{figure*}

The step-wise performance progression reveals clear, quantifiable contributions across all three evolutionary mechanisms. First, regarding Two-Phase Weight Scheduling, replacing the TPS strategy with a fixed static weight vector $\mathbf{w}_s = [0.90, 0.10]$ results in an immediate accuracy drop of 0.86\% and expands the selected subset from 87.8 to 121.7 features, while increasing runtime from 51.3 to 53.6 seconds. The dynamical necessity of TPS is vividly corroborated by the convergence trajectories in Fig.~\ref{fig:ablation_convergence}. During Phase 1 ($t \le 70$), the aggressive feature-size penalty $\mathbf{w}_1 = [0.70, 0.30]$ enforces rigorous dimensional pruning, during which KDS-EL maintains accuracy parity with KDS-EL w/o TPS while selecting substantially fewer features. Exactly upon crossing the scheduled boundary at $t=70$, switching to accuracy-prioritized weights $\mathbf{w}_2 = [0.95, 0.05]$ produces a distinct secondary accuracy acceleration visible in the overall mean trajectory of Fig.~\ref{fig:ablation_convergence}(a) and in challenging benchmarks such as \textit{BASEHOCK}, \textit{COIL20}, \textit{T1D}, and \textit{Madelon}. In contrast, the static variant KDS-EL w/o TPS plateaus prematurely after generation 30, failing to achieve late-stage boundary refinement.

Second, regarding Knowledge-Guided Elite Refinement and Exponential Memory Decay, stripping the KER operator along with its exponential memory decay (reducing KDS-EL w/o TPS to KDS-EL w/o TPS\&KER) inflicts an additional accuracy loss of 1.63\% (falling to 0.9325), while more than doubling the retained feature count from 121.7 to 290.5. As evidenced across all panels of Fig.~\ref{fig:ablation_convergence}, stripping KER and historical feedback depresses the convergence envelope throughout the entire evolutionary timeline, demonstrating that closed-loop feature-importance accumulation is essential for steering elite mutations toward productive dimensions. Third, regarding Mutual Information Initialization and Adaptive Re-initialization, stripping the domain prior seeding and diversity recovery block (MII + AIR)---reducing the framework to the unguided multi-task swarm MEL~\cite{mel2024} initialized uniformly at random with bidirectional weight updates---triggers a severe performance drop: classification accuracy falls by 4.58\% to 0.8868, and the mean feature count surges to 2{,}722.2. The convergence curves in Fig.~\ref{fig:ablation_convergence} confirm that unguided random initialization wastes dozens of generations wandering in dense, non-informative subspaces before making meaningful progress. This substantial 4.58\% performance deficit highlights the foundational role of domain prior guidance in circumventing high-dimensional cold-start blindness, while reflecting the cumulative architectural transition from unguided random exploration to knowledge-seeded dual-subpopulation dynamics.

Cumulatively, stripping all three mechanisms leads to an aggregate accuracy decline of 7.06\% and a 31-fold surge in feature subset size, from 87.8 to 2{,}722.2 features. This confirms that the three mechanisms operate orthogonally and in a complementary manner: MI initialization provides foundational dimensionality reduction, KER delivers fine-grained elite refinement, and TPS harmonizes search objectives across evolutionary stages. Furthermore, the structural design of greedy 1-to-1 in-place elite replacement prevents excessive selection pressure that would otherwise collapse swarm diversity, ensuring robust convergence across heterogeneous problem domains.

\subsubsection{Computational efficiency and runtime analysis}\label{sec:efficiency}

Computational efficiency is a critical evaluation criterion for high-dimensional feature selection, particularly in wrapper-based metaheuristics where repeated classifier evaluations incur substantial runtime overhead. Fig.~\ref{F11} visualizes the overall performance trade-off across classification accuracy, the number of selected features, and computational runtime across the 20 benchmark datasets. The complete per-dataset runtime measurements for all 12 evaluated methods are documented in Supplementary Table S5.

\begin{figure}[!htb]
\centering
\includegraphics[width=0.88\linewidth]{F11_runtime_scatter.pdf}
\caption{Overall performance trade-off across classification accuracy, selected feature subset size, and computational runtime.}
\label{F11}
\end{figure}

As shown in Fig.~\ref{F11} and Supplementary Table S5, KDS-EL achieves an average execution time of 51.3 seconds across the 20 benchmark datasets, outperforming 10 of the 11 comparison methods. Compared with the foundational precursor MEL algorithm, which requires 175.0 seconds on average, KDS-EL delivers a 3.4-fold acceleration. Furthermore, recent top-tier single-objective competitors incur substantial computational penalties: bIRIME requires 326.2 seconds on average (6.4× slower than KDS-EL), while BinCOA requires 433.5 seconds on average (8.5× slower than KDS-EL). On the large-sample 4,862-dimensional \textit{BASEHOCK} benchmark, bIRIME (2,379.5s) and BinCOA (2,362.7s) require nearly 40 minutes per single run due to multi-operator candidate evaluations ($3N$ to $4N$ classifier calls per generation), whereas KDS-EL completes in merely 130.1 seconds, achieving an 18.3-fold wall-clock acceleration. The computational scalability of KDS-EL is also evident on ultra-high-dimensional datasets: on the 54,675-dimensional \textit{T1D} dataset, KDS-EL completes in 118.8 seconds, running more than twice as fast as PSO-CSM at 256.6 seconds and SaWDE at 216.1 seconds, while running nearly 8 times faster than BinCOA at 902.2 seconds. Although the evolutionary filter method DENCA-SO records the fastest runtime (2.1 seconds) because its algebraic evaluation completely avoids inductive classifier calls, it suffers from severe feature redundancy (5,323.4 features) and an 8.62 percentage point accuracy penalty. Similarly, while FWPSO exhibits a marginally faster wrapper runtime of 49.7 seconds due to early-terminating simulated annealing heuristics, it incurs a substantial 3.97 percentage point accuracy deficit and selects over 12 times more features than KDS-EL.

The superior computational scalability of KDS-EL is primarily driven by three synergistic factors. First, in $k$-NN wrapper evaluations, distance computation complexity scales with the number of selected features. By aggressively reducing the active feature count from thousands to dozens during initialization and Phase 1, the per-evaluation cost of the classifier decreases by orders of magnitude. Second, because KER operates exclusively on a compact elite pool with a sparse perturbation fraction, its incremental complexity scales linearly with dimension and remains negligible compared to full population evaluations. Third, unlike complex multi-candidate trial meta-heuristics that invoke three to four classification evaluations per individual per generation, KDS-EL restricts exploratory perturbations to a strictly bounded elite pool ($k_e = 3 \ll N$), keeping generational evaluations linear in population size while accelerating fitness evaluations via extreme feature sparsity. Consequently, KDS-EL establishes an ideal balance of top-tier classification accuracy, extreme feature sparsity, and high computational efficiency.

\subsubsection{Parameter sensitivity and robustness analysis}\label{sec:parametric}

To evaluate the algorithmic robustness of KDS-EL and justify its default parameter configurations, a one-factor-at-a-time parameter sensitivity analysis is conducted following the standard methodology in evolutionary computation~\cite{mimpode2025}. The evaluation is performed across four representative benchmark datasets spanning distinct dimensionalities and sample scales: \textit{Adenoma}, \textit{Colon}, \textit{ALL4}, and \textit{BASEHOCK}. Five key hyperparameters are systematically investigated: the knowledge weight decay factor $\gamma$, the phase split ratio $\kappa$, the elite perturbation fraction $\phi$, the mutual information initialization ratio $\rho$, and the elite pool size $k_e$, while all other parameters are held at their default values in Table~\ref{T3}. Table~\ref{T8} summarizes the tested parameter ranges, recommended optimal configurations, stable operating bands, and maximum accuracy deviations across the four representative datasets. Furthermore, Fig.~\ref{fig:param_sensitivity} illustrates the corresponding accuracy and feature-count trajectories alongside stable operating bands across all five parameters, with complete per-dataset matrices reported in Supplementary Table S6.

\begin{table}[!htb]
\centering
\caption{Summary of one-factor-at-a-time (OFAT) parameter sensitivity analysis across four representative benchmark datasets (\textit{Adenoma}, \textit{Colon}, \textit{ALL4}, and \textit{BASEHOCK}).}
\label{T8}
\small
\setlength{\tabcolsep}{4pt}
\renewcommand{\arraystretch}{1.15}
\begin{tabular*}{\linewidth}{@{\extracolsep{\fill}}lcccc@{}}
\toprule
Parameter & Tested range & Optimal & Stable range & Max $\Delta\mathrm{Acc}$ (pp) \\
\midrule
Knowledge decay ($\gamma$) & 0.90--1.00 & 0.99 & [0.95, 0.995] & $\pm$1.2 \\
Phase split ratio ($\alpha$) & 0.30--0.90 & 0.70 & [0.60, 0.80] & $\pm$0.7 \\
Elite flip fraction ($\phi$) & 0.005--0.10 & 0.01 & [0.005, 0.02] & $\pm$1.3 \\
MI initialization ratio ($\rho$) & 0.0--1.0 & 0.50 & [0.25, 0.75] & $\pm$1.2 \\
Elite pool size ($k_e$) & 1--10 & 3 & [2, 5] & $\pm$0.9 \\
\bottomrule
\end{tabular*}
\end{table}


\begin{figure*}[!htb]
\centering
\includegraphics[width=\linewidth]{F14_parameter_sensitivity.pdf}
\caption{Parameter sensitivity analysis of KDS-EL across five key hyperparameters.}
\label{fig:param_sensitivity}
\end{figure*}

The sensitivity trajectories across the five hyperparameters yield critical insights into the search dynamics and stability of KDS-EL. First, regarding the knowledge weight decay factor $\gamma$, setting $\gamma=0.99$ yields optimal performance with an average accuracy of 0.944 and 107 selected features. An overly low decay factor such as $\gamma=0.90$ causes premature forgetting of valuable feature associations discovered in earlier iterations, degrading accuracy by 1.2\%, whereas disabling decay entirely at $\gamma=1.00$ accumulates obsolete historical feedback and inflates the feature count to 151 without accuracy gain. Second, regarding the phase split ratio $\kappa$, setting $\kappa < 0.50$ allocates insufficient iterations to Phase 1, resulting in inadequate feature pruning with 146 to 189 retained features, while an overly prolonged Phase 1 with $\kappa=0.90$ compresses the Phase 2 refinement window, leading to a 0.7\% accuracy decline. The broad range between 0.60 and 0.80 provides a stable operating regime with accuracy variations remaining strictly below 0.5\%, confirming $\kappa=0.70$ as an ideal balance. Third, regarding the elite perturbation fraction $\phi$, setting $\phi=0.01$ maintains an optimal balance between localized exploitation and structural exploration across heterogeneous dimensions (spanning from 5 flipped features on \textit{madelon} to 546 on \textit{T1D}). Because the greedy acceptance criterion \eqref{eq:inplace} rigorously discards non-improving mutations, $\phi=0.01$ provides sufficient exploratory momentum without provoking excessive trial rejections. Conversely, increasing $\phi \ge 0.05$ introduces excessive perturbation that overwhelms the greedy acceptance filter, resulting in near-zero mutant admission rates, reducing classification accuracy by 1.3\%, and inflating the selected feature count to 203. Fourth, regarding the mutual information initialization ratio $\rho$, disabling MI seeding entirely at $\rho=0.0$ imposes a 1.2\% accuracy penalty and increases execution time by 21\% to 62.3 seconds, while setting $\rho=1.0$ homogenizes the initial population and induces premature convergence. The default balanced setting $\rho=0.50$ effectively reconciles prior domain knowledge guidance with stochastic diversity. Finally, regarding the elite pool size $k_e$, refining only a single elite with $k_e=1$ restricts knowledge propagation across subpopulations, while setting $k_e=10$ dilutes the quality signal of matrix $\mathbf{W}$, establishing $k_e=3$ as the optimal default.

Across all five hyperparameters, the performance trajectories in Fig.~\ref{fig:param_sensitivity} demonstrate wide, forgiving basins of attraction rather than sharp, brittle peaks. The maximum accuracy variation across the identified stable operating bands remains strictly within 1.5\%, and KDS-EL consistently outperforms the comparison algorithms across broad parameter ranges. This demonstrates that KDS-EL exhibits exceptional hyperparameter robustness and plug-and-play adaptability, delivering reliable optimization performance across diverse high-dimensional applications without requiring tedious, problem-specific hyperparameter tuning.

\subsubsection{Discussion}\label{sec:discussion}\label{sec:ceiling}

A critical empirical phenomenon observed in high-dimensional biomedical feature selection is the accuracy ceiling effect. Among the 20 benchmark datasets, eight gene-expression microarrays, including \textit{ALL1}, \textit{ALL\_AML}, \textit{Adenoma}, \textit{DLBCL}, \textit{Gastric2}, \textit{Leukaemia}, \textit{Lymphoma}, and \textit{SRBCT}, exhibit classification accuracy exceeding 0.998 across top-performing methods. Under such near-perfect classification conditions, the few available test instances are completely separated, and accuracy differences among leading algorithms become narrower than 0.002, corresponding to stochastic variations in at most one cross-validation sample. Consequently, classification accuracy alone ceases to provide discriminative utility, and the number of selected features becomes the decisive criterion for evaluating algorithmic efficacy.

Under this saturated regime, KDS-EL demonstrates substantial superiority in feature subset compactness, identifying ultra-sparse biomarker signatures of only 39 to 83 features across all eight ceiling benchmarks. Specifically, KDS-EL extracts minimal subsets of 39 features on \textit{ALL1}, 45 on \textit{Lymphoma}, and 64 on \textit{DLBCL}. In stark contrast, single-objective competitors retain extensive subsets containing hundreds to thousands of redundant features: FWPSO requires 278 to 1{,}362 features to achieve comparable accuracy, BinCOA selects 1{,}136 to 11,041 features (e.g., 6,120 features on \textit{ALL1} and 11,041 on \textit{Gastric2}), evolutionary filter DENCA-SO retains 1,056 to 10,971 features (5,877 on \textit{ALL1}), and the precursor MEL retains 2,483 to 4,248 features. This demonstrates that conventional evolutionary algorithms suffer from severe feature bloat in ceiling scenarios, whereas KDS-EL effectively isolates the true minimal discriminative subset.

This dimensional contrast highlights the profound practical and clinical relevance of KDS-EL. In translational bioinformatics and cancer diagnostics, an ultra-compact signature of 40 to 80 biomarkers can be readily validated through targeted RT-qPCR assays, integrated into biological pathway enrichment analyses, and engineered into cost-effective diagnostic gene panels. In contrast, candidate signatures containing hundreds or thousands of uncurated genes are cost-prohibitive to validate clinically and introduce high risks of false discoveries driven by spurious co-expression artifacts. By reconciling top-tier classification accuracy, extreme biomarker sparsity, and computational scalability, KDS-EL establishes a rigorous and deployable performance frontier for high-dimensional biomedical data analytics.

Beyond nearest-neighbor metrics, a critical concern in wrapper-based evolutionary feature selection is inductive bias coupling: whether the optimized subsets merely fit the local geometric neighborhood of 3-NN while suffering performance collapse when transferred to alternative learning paradigms. To examine classifier-agnostic generalizability, the compact feature subsets discovered by KDS-EL were evaluated under Support Vector Machines with radial basis function kernels and Random Forest classifiers across representative benchmark datasets spanning distinct dimensional scales. As documented in Supplementary Table S9, the compact subsets deliver substantial classification improvements over the unreduced full-dimensional feature sets across all evaluated learning architectures. On the 12{,}625-dimensional \textit{ALL4} dataset, reducing the active feature space by 99.0\% from 12{,}625 to 124 features elevates Random Forest classification accuracy from 78.36\% to 90.35\% and 3-NN accuracy from 69.94\% to 87.25\%, while SVM accuracy reaches 85.96\%. Similarly, on the \textit{Adenoma} benchmark, all three classifiers attain 100.00\% classification accuracy on the 66 selected features. This cross-paradigm consistency confirms that KDS-EL extracts intrinsically salient features rather than overfitting to classifier-specific geometric manifolds.

Furthermore, translating computational feature selection into clinical insight requires that the isolated dimensions reflect genuine pathophysiological mechanisms rather than statistical noise. Granular inspection of the minimal signatures identified by KDS-EL reveals strong concordance with established cancer biology. In the \textit{Colon} adenocarcinoma microarray, top-ranked selected features map directly to well-characterized oncological biomarkers, as detailed in Supplementary Table S10. Specifically, KDS-EL identifies \textit{COL11A1} (probe \textit{Hsa.36689}), a recognized driver of desmoplastic stroma and tumor invasion; \textit{GUCA2A} (Guanylin, probe \textit{Hsa.831}), a mucosal differentiation hormone whose loss constitutes a universal clinical hallmark of early colorectal adenoma progression~\cite{alon1999}; and smooth muscle structural markers \textit{DES} (Desmin, probe \textit{Hsa.3331}) and \textit{ACTG2} (probe \textit{Hsa.1832}), which capture muscularis mucosae integrity versus neoplastic infiltration. In \textit{Prostate} cancer, KDS-EL consistently isolates \textit{HPN} (Hepsin, probe \textit{1005\_at}) and \textit{AMACR} (Alpha-Methylacyl-CoA Racemase, probe \textit{37639\_at}), both clinical immunohistochemical gold standards in prostatic carcinoma diagnosis~\cite{singh2002,rubin2002}. This rigorous biological alignment corroborates that KDS-EL captures genuine molecular drivers of malignancy, bridging evolutionary machine learning with clinical translational oncology.

